\paragraph{Changing who acted on whom.}
If a pictured sentence says that Qin surrounded Chu, exchanging the two names should exchange the besieger and the besieged. We tested this prediction in classical Chinese clauses from a parallel corpus. Two writers received equal exposures to 203 training clauses. One saw each clause twice. The other saw the original and a version with the participants exchanged in both the source and the modern gloss. On twelve held-out event pairs, the exchange-trained writer produced faithful modern answers in both directions for four pairs. Replaying the originals achieved two pairs, as did the parent and the phrase composer (Fig.~\ref{fig:event-roles}). The advantage was specific to these events. On 1,024 further passages, exchange training made 19,057 glyph edits, compared with 19,077 after replay, 18,202 for the parent and 15,564 for the composer. A masked reading accepted none of twelve diverse translations from either trained writer and ten archive references.

We then replaced both state names with historical images wherever both forms were available. The writer had to preserve the supplied strokes and modernize the surrounding sentence. It reproduced both names exactly once in the right order and changed both answers on four of twenty pairs. A reversible assignment of unfamiliar ink to familiar participant shapes increased that count to fourteen. Assignment by overall participant frequency reached thirteen. Both assignment methods produced faithful answers in both directions on two of eight reviewed pairs. In a controlled set holding the bronze forms of Qin and Chu fixed, changing the verb changed the relation expressed for all five verbs in both directions. Six of the ten answers were acceptable complete translations. The rescue answers repeated the rescuer and the destruction answers used an awkward transitive verb. Four of the ten underlying printed prompts had appeared in training. These results establish limited transfer of participant roles to new images in familiar constructions. They leave the reading of an unfamiliar inscription and the generation of dependable book prose open.
